% before-body.tex - Configuración del cuerpo del documento
\newtheorem{teorema}{Teorema}[chapter]
\newtheorem{definicion}{Definición}[chapter]
\newtheorem{lema}{Lema}[chapter]
\newtheorem{corolario}{Corolario}[chapter]
\newtheorem{proposicion}{Proposición}[chapter]


\begin{titlepage}
\hspace{-1.9cm} %Este comando es para mandar a la izquierda ;)
%Aquí empiezan los tres minipages de antes
\begin{minipage}[t][0.03\textheight][c]{0.22\textwidth}
        \includegraphics[scale= 0.23]{UNACH.png}
\end{minipage}\hspace{0.9cm}
\begin{minipage}[t][0.03\textheight][c]{0.69\textwidth}
\begin{center}
                \textsc{\huge Universidad Autónoma de Chiapas}\\[0.3cm]
                \hrule height 2.5pt
                \vspace{0.2cm}
                \hrule height1pt
                \vspace{0.3cm}
                \textsc{\Large Facultad de Ciencias en Física y Matemáticas}
\end{center}
\end{minipage}\hspace{0.2cm}
\begin{minipage}[t][0.03\textheight][c]{0.2\textwidth}
		\includegraphics[width=2.7cm]{logofcfm.png}
\end{minipage}\\
%%%%%%%%%%%%%%%%%%Aquí comienzan los otros dos minipages del título%%%%%%%%%%%%%%%%%%%%%%%%%%%%%%%%%%%%%%%%%%%%%%%%%%%%%%%%%%%%%%%%%%%%%%%%%%%%%%%%%%%%%%%%%%%%%%
\begin{minipage}[t][0.93\textheight][c]{0.06\textwidth}
\vspace{60pt}
    \begin{center}
        \vrule width1pt height18cm
        \vspace{5mm}
        \vrule width2.5pt height18cm
        \vspace{5mm}
        \vrule width1pt height18cm
   \end{center}
\end{minipage}\hspace{1.3cm} %Esto mueve las letras que tienes ahí
\begin{minipage}[t][0.95\textheight][c]{0.76\textwidth}

            \begin{center}
                {\Large\bfseries Aproximación numérica de la solución para ecuaciones diferenciales estocasticas con saltos }\\[1cm]
                \textsc{\huge \textbf{T\, E\, S\, I\, S}}\\[1.5cm]
                \textsc{\large QUE PARA OBTENER EL GRADO DE:}\\[0.3cm]
                \textbf{\textsc{MAESTRA EN CIENCIAS MATEMÁTICAS}}\\[1.5cm]
                \textsc{\large PRESENTA:}\\[0.3cm]
                \textbf{\textsc{\large {MARA  DOMINGUEZ LIMAS  X180099}}}\\[1.5cm]
                {\large\scshape Director de tesis:\\[0.3cm]
                {\textbf{\large Dr. YOFRE HERNÁN GARCÍA GÓMEZ }}}\\[2.0cm]
                \large{Tuxtla Gutiérrez, Chiapas a \today.}\\[1.5cm]  
                {\large{Apoyado con beca SECIHTI y Proyecto Ciencia Básica y De Frontera CBF 2025 -I-4436}}
                
                

            \end{center}
\end{minipage}
\end{titlepage}

\pagebreak[2]

\chapter*{Dedicatoria}
\begin{center}
\itshape
\setlength{\parindent}{0pt} % Elimina sangría para mayor elegancia
\setlength{\parskip}{1.2\baselineskip} % Espacio generoso entre párrafos


\end{center}


\chapter*{Agradecimientos}


\chapter*{Resumen}

En este trabajo se aborda el problema de la aproximación numérica de
soluciones de ecuaciones diferenciales estocásticas con difusión y
saltos (EDESJs), con aplicación al modelo de Cramér-Lundberg extendido
para el cálculo de reservas técnicas en el sector asegurador de vivienda
de la Ciudad de México. Se establece un marco riguroso de cálculo
estocástico que incluye la demostración de existencia y unicidad de
soluciones fuertes bajo condiciones de Lipschitz local y crecimiento
lineal global; la isometría de Itô para la parte de saltos; y las
desigualdades de Burkholder-Davis-Gundy en espacios de Hilbert.

Se implementan y analizan dos esquemas numéricos: el método de Milstein
compensado, que alcanza convergencia fuerte de orden \(\gamma = 1/2\)
para la componente difusiva, y el esquema Semi-implícito compensado, que
trata el término de interés \(r(t)U(t)\) de forma implícita,
garantizando estabilidad incondicional. Bajo hipótesis de regularidad de
los coeficientes (Lipschitz global, crecimiento lineal, momentos finitos
de saltos), se demuestra que ambos esquemas convergen fuertemente con
orden \(\gamma = 1/2\), utilizando técnicas de descomposición del error
local, estimaciones recursivas y la desigualdad de Grönwall discreta.

La calibración del modelo incorpora datos reales del mercado mexicano:
primas ajustadas por inflación (INEGI), tasa de inversión combinada con
activos locales (CETES) y globales (bonos USD), volatilidad estacional
basada en el índice ILS, e intensidades de siniestros
hidrometeorológicos (CONAGUA) y sísmicos (SSN). Mediante simulación
Monte Carlo (\(N_{\text{sim}} = 1000\)) se comparan los esquemas en
términos de mediana final, intervalos de confianza al \(90\%\) y
detección de saltos, validando la robustez del modelo.

Los resultados muestran que el esquema semi-implícito compensado
presenta ventajas críticas para la gestión de riesgos en la CDMX:

\begin{enumerate}
\def\labelenumi{\roman{enumi}.}
\item
  Prudencia regulatoria (\(+2.30\%\) en reserva mediana);
\item
  Estabilidad incondicional ante shocks en tasas de interés;
\item
  Eficiencia computacional (reducción del 96.7\% en pasos de tiempo para
  horizontes extendidos). Se concluye que, para sistemas dominados por
  saltos catastróficos y volatilidad macroeconómica, el esquema
  semi-implícito es preferible sobre Milstein.
\end{enumerate}

\textbf{Palabras clave}: Ecuaciones diferenciales estocásticas con
saltos, método de Milstein, esquema semi-implícito, convergencia fuerte,
modelo de Cramér-Lundberg, simulación Monte Carlo, gestión de riesgos,
seguros de vivienda, Ciudad de México.

\chapter*{Abstract}

This work addresses the numerical approximation problem of stochastic
differential equations with diffusion and jumps (SDEJs), applied to an
extended Cramér-Lundberg model for calculating technical reserves in the
housing insurance sector in Mexico City. A rigorous stochastic calculus
framework is established, including the demonstration of existence and
uniqueness of strong solutions under local Lipschitz and global linear
growth conditions; Itô isometry for the jump component; and
Burkholder-Davis-Gundy inequalities in Hilbert spaces.

Two numerical schemes are implemented and analyzed: the compensated
Milstein method, which achieves strong convergence of order γ = 1/2 for
the diffusive component, and the compensated semi-implicit scheme, which
treats the interest term r(t)U(t) implicitly, guaranteeing unconditional
stability. Under regularity hypotheses of the coefficients (global
Lipschitz, linear growth, finite jump moments), it is demonstrated that
both schemes converge strongly with order γ = 1/2, using local error
decomposition techniques, recursive estimates, and the discrete Grönwall
inequality.

The model calibration incorporates real data from the Mexican market:
inflation-adjusted premiums (INEGI), combined investment rate with local
(CETES) and global (USD bonds) assets, seasonal volatility based on the
ILS index, and intensities of hydrometeorological (CONAGUA) and seismic
(SSN) claims. Through Monte Carlo simulation (N\_sim = 1000), the
schemes are compared in terms of final median, 90\% confidence
intervals, and jump detection, validating the robustness of the model.

The results show that the compensated semi-implicit scheme presents
critical advantages for risk management in Mexico City:

\begin{enumerate}
\def\labelenumi{\roman{enumi}.}
\item
  Regulatory prudence (+2.30\% in median reserve);
\item
  Unconditional stability against interest rate shocks;
\item
  Computational efficiency (96.7\% reduction in time steps for extended
  horizons). It is concluded that for systems dominated by catastrophic
  jumps and macroeconomic volatility, the semi-implicit scheme is
  preferable over Milstein.
\end{enumerate}

\textbf{Keywords:} Stochastic differential equations with jumps,
Milstein method, semi-implicit scheme, strong convergence,
Cramér-Lundberg model, Monte Carlo simulation, risk management, housing
insurance, Mexico City.